\section{Real-Binary Validation and Gaia DR3/DR4 Tests}
\label{sect:realdata}

\subsection{Legacy GJ~765.2 Visual--SB2 Data}

As a first real-system check of the physical orbit engine, we use a legacy subset for the nearby visual and double-lined spectroscopic binary GJ~765.2 = HD~186922 = HIP~96656. The later combined interferometric--spectroscopic solution of \citet{Balega2007} provides a useful reference: $P=11.919\pm0.003$~yr, $e=0.240\pm0.003$, $i=80.2\pm0.2^\circ$, $\Omega_{\rm asc}=293.0\pm0.6^\circ$, $\omega_{\rm rel}=250.0\pm0.1^\circ$, component masses $0.831\pm0.020$ and $0.763\pm0.019\,M_\odot$, and orbital parallax $31.0\pm0.5$~mas.

The legacy subset used here contains 11 resolved relative positions and 44 paired CORAVEL epochs for each RV component. Most of the astrometry predates the high-precision speckle sequence used in the later reference solution; the final two positions in the subset are 1993 speckle measurements with quoted 5~mas positional uncertainties. Position angle and separation are converted to the tangent plane through
\begin{equation}
\Delta\alpha^*=\rho\sin\theta,\qquad
\Delta\delta=\rho\cos\theta,
\end{equation}
where $\theta$ is measured from North through East. The fourth astrometric column is treated as an isotropic one-sigma tangent-plane positional uncertainty. This is equivalent, to first order, to the polar weighting used by the legacy orbit code. The two RV components retain their quoted measurement uncertainties.

No Gaia epoch measurements are present in this data set. The Gaia likelihood is therefore disabled and $\betag$ is irrelevant. The fit varies the ten physical parameters $P$, $T_0$, $e$, $i$, $\omega_1$, $\Omega$, $M_1$, $M_2$, $\varpi$, and $\gamma$. The header values are used only for initialization, not as priors. Because visual astrometry alone has the usual $180^\circ$ node ambiguity, both node branches are initialized and the lower-$\chi^2$ joint astrometry--SB2 solution is retained.

\subsection{Fit to the Legacy Subset}

The real-data fit contains 110 scalar constraints and ten free parameters. It converges to
\begin{equation}
\chi^2=104.63,\qquad \nu=100,\qquad \chi^2_\nu=1.046.
\end{equation}
Table~\ref{tab:gl765} compares the recovered orbit with the later combined solution. The legacy subset is not statistically independent of \citet{Balega2007}, because the CORAVEL velocities overlap with the data used in that work; this comparison is therefore a consistency and implementation check rather than an independent astrophysical measurement.

\begin{table}
\begin{center}
\caption{GJ~765.2 Legacy-Subset Consistency Check}
\label{tab:gl765}
\begin{tabular}{lcc}
\hline\noalign{\smallskip}
Quantity & RAA legacy-subset fit & \citet{Balega2007}\\
\hline\noalign{\smallskip}
$P$ (yr) & $11.728\pm0.073$ & $11.919\pm0.003$\\
$e$ & $0.2489\pm0.0101$ & $0.240\pm0.003$\\
$i$ (deg) & $81.83\pm1.37$ & $80.2\pm0.2$\\
$\Omega_{\rm asc}$ (deg) & $289.07\pm3.27$ & $293.0\pm0.6$\\
$\omega_{\rm rel}$ (deg) & $251.79\pm2.25$ & $250.0\pm0.1$\\
$M_1+M_2$ ($M_\odot$) & $1.5897$ & $1.594$\\
$\varpi_{\rm orb}$ (mas) & $35.44\pm2.24$ & $31.0\pm0.5$\\
\noalign{\smallskip}\hline
\end{tabular}
\end{center}
\tablecomments{0.94\textwidth}{The quoted uncertainties for the RAA fit are local formal one-sigma errors from the weighted least-squares Jacobian. Individual component masses are not compared one-to-one because the component labels in the legacy RV file do not map unambiguously onto the later A/B convention. The total mass is label invariant.}
\end{table}

The total mass differs from the later combined value by only $0.27\%$. The orbital parallax is less precise and remains higher than the later value, but is within about two combined formal standard deviations. More importantly, the legacy header parallax of 54.27~mas is strongly disfavoured by the joint data: fixing $\varpi=54.27$~mas increases the fit to $\chi^2=168.94$ ($\chi^2_\nu=1.673$ for 101 degrees of freedom), whereas fixing the published $31.0$~mas gives $\chi^2=108.70$ ($\chi^2_\nu=1.076$). Thus the resolved orbit and the two RV curves correctly force the dynamical scale away from the inconsistent legacy starting parallax.

This experiment validates only the Newtonian/SB2 plus resolved-relative-astrometry core on a real stellar system. It does not validate the marginal-resolution Gaia response, because no Gaia epoch along-scan coordinates or images are included.

\subsection{Gaia DR3 as an Intermediate Catalogue-Level Test}
\label{sect:dr3bridge}

Gaia DR3 cannot provide the general source-level epoch along-scan astrometry needed to refit the response hierarchy $\mzero$--$\mtwo$ directly, but it does provide two useful catalogue-level diagnostics. First, the DR3 \texttt{gaia\_source} table publishes image-parameter-determination (IPD) quantities that retain information about partial resolution. The field \texttt{ipd\_frac\_multi\_peak} is the percentage of successful IPD windows in which more than one peak was identified, whereas \texttt{ipd\_gof\_harmonic\_amplitude} measures the amplitude of scan-angle-dependent variation in the IPD goodness of fit.\footnote{Gaia DR3 \texttt{gaia\_source} data model: \url{https://gea.esac.esa.int/archive/documentation/GDR3/Gaia_archive/chap_datamodel/sec_dm_main_source_catalogue/ssec_dm_gaia_source.html}, accessed 2026 August 17.} The accompanying harmonic phase is defined modulo $180^\circ$. The Gaia documentation notes that for binaries with separation $\lesssim0.1$~arcsec the phase can trace the binary position angle when the profile remains effectively single-peaked, while more clearly resolved doubles can instead have a large multi-peak fraction. These quantities are therefore relevant empirical diagnostics of the same scan-orientation dependence explored by the resolution-aware surrogate, although they are not direct measurements of its Gaussian mode-splitting threshold.

Second, the DR3 \texttt{nss\_two\_body\_orbit} table publishes non-single-star orbital solutions. For astrometric orbital models, the fitted Thiele--Innes constants $A$, $B$, $F$, and $G$ describe the orbit of the photocentre and can be transformed to the Campbell elements $a_0$, $i$, $\omega$, and $\Omega$.\footnote{Gaia DR3 \texttt{nss\_two\_body\_orbit} data model: \url{https://gea.esac.esa.int/archive/documentation/GDR3/Gaia_archive/chap_datamodel/sec_dm_non--single_stars_tables/ssec_dm_nss_two_body_orbit.html}, accessed 2026 August 17.} If GJ~765.2 has a usable astrometric NSS solution, this provides a direct catalogue-level comparison between the Gaia photocentre orbit and the independently constrained visual--SB2 orbit.

For that comparison, the ordinary unresolved-photocentre model predicts
\begin{equation}
 a_{\rm phot}=|\beta-B|a_{\rm rel},\qquad
 B=\frac{M_2}{M_1+M_2}.
\end{equation}
Using the published GJ~765.2 masses gives $B=0.47867$. The published component $V$-magnitude difference, $\Delta V=0.65$~mag, gives the temporary optical light-fraction proxy
\begin{equation}
\beta_V=\frac{1}{1+10^{0.4\Delta V}}=0.35465.
\end{equation}
With $a_{\rm rel}=189$~mas, the corresponding $\mzero$ benchmark is
\begin{equation}
 a_{\rm phot,V}=|0.35465-0.47867|\times189
 =23.44~{\rm mas}.
\label{eq:gl765_phot_benchmark}
\end{equation}
The subscript $V$ is essential: this is not a measured Gaia $G$-band light fraction and is used only to define a transparent external benchmark.

To avoid a potentially unsafe coordinate-only association for this high-proper-motion binary, the target-specific DR3 query implemented with the analysis code identifies HIP~96656 through the Gaia Hipparcos-2 best-neighbour cross-match and then joins the matched source to \texttt{gaia\_source} and \texttt{nss\_two\_body\_orbit}. The cross-match table links each Hipparcos-2 object to its best Gaia neighbour and records whether the Gaia position was propagated with the available astrometric parameters.\footnote{Gaia DR3 archive table inventory includes \texttt{hipparcos2\_best\_neighbour}; see \url{https://gea.esac.esa.int/archive/documentation/GDR3/}.} The query, parser, Thiele--Innes conversion, and regression tests are versioned with the code. At the time of this manuscript update, no target-specific DR3 numerical values are included in the scientific results because a validated catalogue export for HIP~96656 has not yet been ingested. Accordingly, DR3 is treated here as an intermediate catalogue/IPD validation layer rather than as evidence for or against any one response model.

The intended DR3 comparison is therefore deliberately limited. A recovered Gaia astrometric NSS orbit can be compared with the externally known $P$, $e$, $i$, and with the $23.44$~mas ordinary-photocentre scale in Equation~(\ref{eq:gl765_phot_benchmark}); RUWE, excess-noise and IPD diagnostics can test whether Gaia reports the expected symptoms of scan-dependent source structure. None of these quantities permits a likelihood-level comparison of $\mzero$, $\mone$, and $\mtwo$ because the individual stellar AL measurements used by the DR3 astrometric solution are not generally available.

\subsection{Official DR4 Products and GJ~765.2 as a Target}

The official Gaia DR4 expected-content page lists the two products most directly relevant to the next validation step: \texttt{epoch\_astrometry}, containing individual astrometric measurements for sources treated in the main astrometric processing, and \texttt{epoch\_image}, containing pre-processed sky-projected CCD sample values. It also lists \texttt{rvs\_epoch\_data\_double} for double-lined RVS transits.\footnote{Gaia DR4 expected content: \url{https://www.cosmos.esa.int/web/gaia/dr4}, accessed 2026 August 17.} The release is based on 66 months of observations, from 2014 July 25 10:30 UTC to 2020 January 20 22:00 UTC, with reference epoch J2017.5; the expected total release volume is approximately 500~TB, dominated by lower-level products. The same official page states that the content remains under development and can change before release. The DR4 overview still labels the release as ``Coming up'' at the time of writing,\footnote{Gaia DR4 overview: \url{https://www.cosmos.esa.int/web/gaia/data-release-4}, accessed 2026 August 17.} while the official release-scenario page gives the target only as ``not before mid 2026''.\footnote{Gaia data-release scenario: \url{https://www.cosmos.esa.int/web/gaia/release}, accessed 2026 August 17.}

GJ~765.2 is a useful prospective DR4 target because its published orbit naturally crosses a wide angular-separation range during the 66-month DR4 interval. Propagating the \citet{Balega2007} orbit gives a sky-projected component separation from approximately 24.7 to 199.1~mas over that interval. The published absolute $V$ magnitudes differ by 0.65~mag, corresponding to the temporary $V$-band light-fraction proxy in Equation~(\ref{eq:gl765_phot_benchmark}). This is not assumed to equal the Gaia $G$-band fraction; it is used only for a pre-release feasibility calculation.

For the current equal-width research baseline $\alpha=50$~mas, the exact equal-width surrogate mode-splitting criterion at $\beta=\beta_V$ gives a critical projected along-scan separation of approximately 128.5~mas. A uniform grid in time and scan orientation over the DR4 interval places 24.1\% of the time--orientation combinations in the surrogate multi-peak regime. For 58.0\% of the interval the instantaneous sky separation is large enough that at least some scan orientations can be multi-peaked; at the most favourable epoch, 55.3\% of possible scan orientations cross the threshold.

These percentages are an orientation envelope, not Gaia transit statistics. They do not use the source's actual accepted DR4 observations, and the 50~mas width is not a calibrated Gaia PLSF parameter. The calculation establishes only that GJ~765.2 spans the transition regime targeted by this work. Once DR4 is public, the correct test is to replace the uniform-angle envelope with the released epoch measurements and, where the profile is not representable by a unique coordinate, to move directly to an \texttt{epoch\_image} sample-domain likelihood.
